A growing fraction of web text is generated by large language models (LLMs).
If future foundation models train on this text, single-model recursive
training is known to cause distributional collapse. We ask the
ecosystem-level analogue: is there a \emph{minimum viable population} (\mvp)
of LLMs that, sharing a common synthetic-text substrate, avoids collapse?
And what kind of inter-agent differentiation actually counts? We run two
complementary harnesses on the same metric suite. In a fine-tuning
ecosystem (\distilgpt copies retraining on each others' \wikitext outputs)
even $N{=}16$ agents collapse, with perplexity rising $2.10\times$ over
six iterations and a slope of $16$~ppl/iter that is statistically
significant ($p < 10^{-6}$); the rate of collapse falls $\sim 5\times$
from $N{=}1$ to $N{=}16$ but never reaches zero. In a retrieval-augmented
ecosystem (frozen API LLMs sharing a growing post pool), \emph{$N{=}2$
LLMs from different pretraining families} preserve full inter-agent
diversity over $T{=}12$ iterations, with both pairwise distance and the
Vendi score yielding regression slopes indistinguishable from zero
($p \approx 0.4$). Holding $N{=}4$ fixed and varying the axis of
differentiation, four agents from different pretrained families preserve
$\mathbf{1.9\times}$ the inter-agent semantic distance of four agents
varied by prompt persona, by retrieval slice, or by random seed alone.
Together these results sharpen the population view of model collapse: the
training-time link is what matters, and the only individuality axis that
buys real collapse-resistance is diversity in the prior, not the
conditioning.
